\section{Identifying the cap moments at distinct sites}\label{sec:site-moments}

We compare the scalar trace with a contraction of the two disk spin caps when each row argument is a distinct site parameter. The comparison first determines the unnormalized scalar factor. Its empty-row case then fixes the denominator used for arbitrary repeated-row moments.

Here and below an integer shift in a trigonometric argument denotes that integer times \(l\); brackets mean twice the sine. Use \(d,e,H,\Psi,Z\) from Sections~\ref{sec:local} and~\ref{sec:vacuum} and \(W,S,R_{\rm pair},V_{\rm pair},G,D_i,L,d_{\rm full}\) as in Lemma~\ref{lem:residue} and Proposition~\ref{prop:common-twist}. Thus \(R_{\rm pair},V_{\rm pair}\) here are the \emph{scalar} functions called \(R,V\) there, and \(d_{\rm full}(w)=\prod_b[w+3-b][w+2-b]\). Write \(\tau_N\) for the trace \(\tau\) on \(N\) sites and roots. We work first with ordered sites \(\mathbf b\) in distinct shift orbits and generic \(K\ne0\). A diagonal trace means the holomorphic specialization of the rational independent-twist trace near a generic diagonal \(K_i=K\), as in Proposition~\ref{prop:common-twist}; it includes contributions from colliding roots.

For this \(N\)-site spin space write \(\Psi=\Psi_N\). Define \(\Lambda\) by taking the transpose of the right vacuum for the reversed site list, reversing its tensor slots back to the original order, and negating all spins. Since \(R^t\) equals simultaneous spin-negation conjugation of \(R\), this is the left vacuum at twist \(-i\). Reversal and negation together preserve matching tensors, so \(\Lambda\) lies in \(D^t\), has empty coefficient one, and is generically unique with that normalization. Moreover \(\Lambda\Psi=1\): two matching tensors pair only on identical occupancy, and an innermost arc then contracts the other matching by the pairing rules, eventually closing a zero loop unless the occupancy was empty.

For a subset \(X\) of size \(k\), with complement \(Y\) of size \(m=N-k\), the comparison before normalizing the scalar trace is
\begin{equation}\label{eq:site-moment-unnormalized}
 \Lambda \prod_{x\in X}T^{K^{-1}}(x)\Psi
 =\frac{(2\sqrt2)^N}{N!\,Z_N^2}
  \prod_{\{b,b'\}\subset\mathbf b}H(b-b')H(b'-b)\,
  \tau_N\!\left(\prod_{x\in X}\frac{L(x)}{d_{\rm full}(x)}\right).
\end{equation}
The product over \(\{b,b'\}\) takes each unordered pair of distinct sites once. The empty subset and \(\Lambda\Psi=1\) give the exact normalization
\begin{equation}\label{eq:trace-normalization}
 \tau_N(1)=
 \frac{N!\,Z_N^2}
 {(2\sqrt2)^N\displaystyle\prod_{\{b,b'\}\subset\mathbf b}
                     H(b-b')H(b'-b)}.
\end{equation}
This value is independent of the twist, including the independent twists before specialization. For \(N=0\), there is one empty root fiber, with \(W=J=1\); with \(Z_0=Z_1=1\), \(0!=1\), and empty pair products equal to one, the formula gives \(\tau_0(1)=1\) and \(\tau_1(1)=1/(2\sqrt2)\). In particular \(\tau_N(1)\ne0\) generically. Substituting this value in \eqref{eq:site-moment-unnormalized} gives
\begin{equation}\label{eq:site-moment}
 \Lambda \prod_{x\in X}T^{K^{-1}}(x)\Psi
 = \frac{\tau_N\!\left(\prod_{x\in X}L(x)/d_{\rm full}(x)\right)}{\tau_N(1)}.
\end{equation}
These identities extend as identities of meromorphic expressions in the site and twist parameters, with pointwise evaluation at regular values.

We use two further consequences of the cap construction. By the exchange and rotation identities of Section~\ref{sec:local}, \(\Lambda\) absorbs adjacent spectral swaps at exchanged arguments (the \(\check R\) are symmetric and preserve empty coefficient on matching tensors), and twisted rotations with twist \(-i\). On regular fusion divisors \(\Lambda E_r=\Lambda_{\rm red}\): uniqueness applies on the reduced list, the pairing rules give \(E_r^tD\subset D\) by contracting adjacent slots via \(C^t\) and squeezing slots with a single spin, and the empty normalization is preserved. The simultaneous disjoint specializations used here and for deletion fusions follow as well, since the symmetric \(Z\) satisfies the fusion recurrences with nonzero denominator at generic independent centers.

Permute the listed sites to the front by exchange. At the first site itself the transfer equals a left rotation weighted by \(K^{-s_1}\) followed by adjacent swaps taking that rapidity back to the front (put \(R_{a1}=P_{a1}\) in the trace). Order the transfers with this evaluation rightmost by commutativity. Push those swaps from the rightmost transfer past remaining transfers into \(\Lambda\); now evaluate the next transfer at the new first site and repeat. Absorbing the accumulated rotations at twist \(-i\) gives \(\Lambda\eta^{\sum_{x\in X}s_x}\Psi\), \(\eta=i/K\).

\subsection{Resolving the finite spin words}

For an ordered block $I=(i_1,\ldots,i_r)$, define its diagonal auxiliary entries by
\[
 A^{I}_{\pm}(w)=\langle\pm|_a
 R_{a i_1}(w-b_{i_1})\cdots R_{a i_r}(w-b_{i_r})|\pm\rangle_a.
\]
Thus $A_+(w)$ below means $A^X_+(w)$ and $A_-(w)$ means $A^Y_-(w)$. We resolve the cap contraction by inserting the commuting spectral projectors of these auxiliary diagonal entries in the partial ordered row monodromies. Entries with a fixed extremal auxiliary spin commute at different arguments (fully polarized Yang--Baxter component, both as row and column). Within each block they conserve total spin and are triangular (rightmost changed label increases), with labels \(\sigma=s\) on \(X\), \(\sigma=-s\) on \(Y\); diagonal signatures are products over its sites of \(d,h,e\) at argument \(w-b\) for \(\sigma=+,0,-\), respectively. These signatures separate simple eigenlines by numerator zeros in each site's shift orbit on using a word-independent common denominator. The term of a joint word vanishes or both block label sums equal some \(r\), with multiplier \(\eta^r\), by conservation. Empty blocks have the obvious scalar conventions. It remains to evaluate \(\Lambda\Pi\Psi\) for each joint projector \(\Pi\).

The next reduction removes zero labels while preserving this contraction. Split each zero site into the ordered pair \(b+1,b-1\), replacing its label by \((+,-)\). The monodromy fusion identity intertwines the block entries through the spin injection \(E_2\), and the signature matches there because \(d(t-l)e(t+l)=h(t)\). The signature remains simple on the fusion locus: in the common pair denominator representation the first numerator choices have zeros (relative to center) \((4,3),(4,1),(2,1)\), the second \((2,1),(2,-1),(0,-1)\) modulo 8. Separation at a generic test argument per block therefore makes the projector regular there. Applying vacuum fusion on both sides proves that \(\Lambda\Pi\Psi\) is unchanged. We can consequently compute the all-nonzero word on an enlarged generic list, performing all exchanges there before restricting the split pairs.

For an all-nonzero word, the extremal orders identify the selected eigenlines. Within each block reorder rapidities to the minimum word (+ then -), or maximum word (- then +), keeping internal same-label orders fixed. Adjacent exchanges transport the projectors via Yang--Baxter. For fixed sum these words are extremal even allowing zeros, so the left eigenline in minimum order is the pure minimum covector, and the right eigenline in maximum order the pure maximum ket (simultaneously for both blocks). If \(U\) is the product of adjacent swaps taking maximum list to minimum, this gives the contraction
\[
 \frac{\Psi(\min)_{\min}\ \Lambda(\max)_{\max}}
 {\langle\min|U|\max\rangle},\qquad
 \langle\min|U|\max\rangle=\prod_{\text{within block }(+,-)} e(b_+-b_-).
\]
Indeed pass each + from first to last in their relative order across the minuses; to arrive still + it could not have donated any label spin, since each encountered minus is initially minimal; induction leaves the minuses unchanged. All these operations precede restriction to fused values.

\subsection{Matching spin words with frozen slots}

The projector calculation has reduced the matrix element to explicit products. We now assign a frozen-slot term to each word and check its phase and denominator. Recall that the label is \(\sigma=s\) on \(X\) and \(\sigma=-s\) on \(Y\), where \(s\) is the actual eigenword spin. Encode an original site's spin by the slot length \(j=1+s\): its slots are \(x-a\) on \(X\) or \(y+a\) on \(Y\), for \(a=1,\ldots,j\). Since both block label sums are \(r\), the numbers of occupied slots are \(k+r\) on the left and \(m-r\) on the right, hence \(N\) in total.

To calculate the word's contribution, let its enlarged plus/minus \(\sigma\)-label counts on \(X,Y\) be \((A,B),(C,D_0)\); in particular a plus label on \(Y\) is an actual minus spin. Conservation gives \(M=A+D_0=B+C\). The minimum ket has actual spins \(+^A-^M+^{D_0}\), and the reversed negated maximum word has actual spins \(+^C-^M+^B\). Rotate the first plus segment to the end in each and use the two-block component formula \eqref{eq:two-block-component}; the phases are \(p^{2A-M+2C-M}=p^{2r}\), converting \(\eta^r\) to \(K^{-r}\). Besides \(K^{-r}\), the result has \(C_0=2\cos 3l\) per enlarged site divided collectively by \(Z_{\rm enlarged}^2\). Within each block equal labels contribute \(H(t)H(-t)\), \(t\) the difference; opposite labels contribute \(1/e(b_+-b_-)\). Across blocks, equal \emph{actual} spins \(n\) give \(H(-n(x-y))^2\); unequal spins give 1.

Applying the fusion recurrence to \(Z_{\rm enlarged}^2/Z_N^2\) reorganizes this contribution as
\[
 \frac{K^{-r} (C_0\alpha)^N}{Z_N^2}\,
 \prod_{\{b,b'\}\subset\mathbf b}H(b-b')H(b'-b)\prod_{x,y}d(x-y)
 \ (-1)^{km+m}\operatorname{Res}_{\mathcal L}
 W\prod_{i,x}S(t_i-x)^{-1}\prod_{i\ {\rm left}}D_i ,
\]
where $\alpha=4\cos l$ and \(\mathcal L\) uses any one labeling of the occupied slots.

Here are details of the elementary reorganization. Besides the base pair products displayed, each occupied slot \(p'\) with another original site \(b\) has a factor \(V_{\rm pair}(p'-b)\) within a block and \([p'-b\mp1]^{-2}\) across (upper sign left slot). Slots at distinct sites interact by \(G(u)\) within block and \(([u][u-2]/[u-1])^2\) across, \(u=\) left minus right there. One verification uses offset/spin lists
\(B_0=((0,-)), B_1=((1,+),(-1,-)), B_2=((0,+))\) (reflect offsets on Y). Explicitly put \(F_{ij}(t)=[t+3]^2[t+4]^2[t+5]^2\) for both indices 1, \([t+4]^2\) for just one, 1 otherwise (squared fusion factors). For \(z=t+a-b\) in the first identity, \(z=t+a+b\) in the second,
\[
\begin{gathered}
 \prod_{(a,n)\in B_i,(b,n')\in B_j}
 \begin{cases}H(z)H(-z)&n=n'\\H(-nz)/([nz-1][nz])&n\ne n'\end{cases}\\
 =F_{ij} H(t)H(-t)\prod_{c=1}^i V_{\rm pair}(t-c)\prod_{h=1}^j V_{\rm pair}(-t-h)\prod_{c,h}G(t-c+h),\\[4pt]
 \prod_{(a,n)\in B_i,(b,n)\in B_j} H(-nz)^2
 =F_{ij} H(t)^2\prod_c[t-c-1]^{-2}\prod_h[t-h-1]^{-2}\\
 {}\times\prod_{c,h}\left(\frac{[t-c-h][t-c-h-2]}{[t-c-h-1]}\right)^2 .
\end{gathered}
\]
Both follow by inserting the three displayed lists, canceling with \([-v]=-[v], [v+8]=-[v]\). Use reflection of all rapidities for within Y. Single-site constants are \(C_0, C_0^2/((2\cos 2l)^2 e(2l))=1,C_0\); equivalently \(C_0\alpha^j\) times the successive left-slot residues of \(\prod V_{\rm pair}\prod G\) at that site alone (residues \(1,1/(2[2l]),[l]^2/(4[2l]^2)\) in ordinary angle coordinates), or their reflected right-slot residues multiplied by \((-1)^j\). Remaining factors match the displayed integrand since across blocks \(V_{\rm pair}S^{\pm1}\) and \(G R_{\rm pair}\) in oriented arguments give the stated cross corrections each with a minus. The total sign including reflection is \((-1)^{(k+r)m+(m-r)k+(m-r)(k+r)+m-r}=(-1)^{km+m}\).

Finally \(K^{-r}=K^k\prod_{i\ {\rm left}} K^{-1}\) and
\(\prod_x L(x)/d_{\rm full}(x)=(-1)^{Nk}K^k\prod_{xy}d(x-y)\prod_{i,x}S(t_i-x)^{-1}\).
Each unlabelled frozen assignment has exactly \(N!\) labelings. At diagonal twists the residue lemma therefore turns the sum of the displayed word contributions into \eqref{eq:site-moment-unnormalized}: its sign agrees because \(N+Nk\equiv km+m\pmod2\), and its prefactor is
\[
 \frac{(C_0\alpha)^N}{N!\,Z_N^2}
 \prod_{\{b,b'\}\subset\mathbf b}H(b-b')H(b'-b),
 \qquad C_0\alpha=(2\cos3l)(4\cos l)=2\sqrt2.
\]
For \(X=\varnothing\), the cap contraction is one, proving \eqref{eq:trace-normalization}. In the independent-twist residue lemma this same empty subset has only right slots and no left twist multipliers. Hence \(\tau_N(1)\) is already independent of every \(K_i\) and equals its generic diagonal value. Substitution of the normalization now proves \eqref{eq:site-moment}.
